\section{Study Design and Methodology}
\subsection{Participants}
Five participants completed the full study. Participants varied in their self-reported spreadsheet expertise, which was measured on a ten-point scale (0=no experience, 10=expert), with a mean expertise of 3.4, indicating a predominantly beginner-to-intermediate knowledge level.
% Insert table of particpants 
\subsection{Tasks}
The Participants completed three spreadsheet comprehension tasks of increasing complexity. Before the participants completed the tasks, they completed a short five minutes tutorial about all the basics of the tool.
\subsubsection{Task 1 - Sales Commissions}
This is a single worksheet in which each employee's final commission is computed through a normal calculation chain, involving VLOOKUP and IF functions. Five questions to this task were asked, about comprehension, prediction and error correction.
\subsubsection{Task 2 - Student Grading}
Again this is a single worksheet but with many IF statements and logic conditions. The six questions of this task aimed to test the understanding of conditional logic. In contrast to the other tasks only one student's data was populated in order to focus even more on the conditional logic.
\subsubsection{Task 3 - Inventory Pricing}
Unlike the other tasks, this workbook contained 3 worksheets with pricing analysis, product catalogue and supplier costs. This task's goal was to test the comprehension when multiple worksheets are involved.

\subsection{Instruments}
Three different instruments were used to collect study data:
\subsubsection{NASA-TLX}
The NASA Task Load Index~\cite{hart1988nasa} is a widely used multidimensional rating procedure for assessing perceived workload. In this study the physical demand was omitted because it is not meaningful for a screen-based study. The five collected subscales were mental demand, temporal demand, self perceived performance, effort and frustration. Each of them had to be rated on a 0-100 scale. After each task the participants had to answer the NASA-TLX without the physical demand, in total 3 times.
%TODO: write more what a good score means and when it is a good score
\subsubsection{Custom UEQ+}  
The study used the User Experience Questionnaire Plus~\cite{Schrepp_Thomaschewski_2019} and the participants had to answer it after they completed all three tasks. The questionnaire consisted of seven scales, Usefulness, Dependability, Efficiency, Perspicuity, Clarity, Intuitive and Visual Aesthetics. Each of these scales had four items on a seven-point Likert scale (1=most negative pole; 7=most positive pole). For the evaluation the scores are transformed to a scale from -3 to +3, where 0 represents a neutral evaluation.
%TODO: explain why the study used exactly these scales
\subsubsection{Open-ended questions}
At the end the participants were invited to answer three open ended questions:
\begin{enumerate}
    \item What they liked most about the tool
    \item What they found frustrating about the tool
    \item What they would improve
\end{enumerate}